\documentclass[11pt,a4paper]{article}
\usepackage{notebooks}

\notebooktitle{Cloud Computing}
\notebooksubtitle{CAT7002 Course Notes}
\notebooktagline{One part per unit. Unit 1: virtualization concepts and hypervisors.}
\notebookauthor{by Bhuvnesh Verma}
\notebooksource{CAT7002 Cloud Computing, RBU Nagpur}

\begin{document}
\notebookcover

\part*{Unit 1: Virtualization Concepts and Hypervisors}

\section{Introduction to Servers}

\field{Short description:}{A server is a computer or software system that provides services,
resources, or data to other computers, called clients, over a network.}

\field{Everyday examples:}{}

\begin{itemize}\itemsep2pt
  \item Opening a website: the browser sends a request to a \textbf{web server}.
  \item Checking Gmail: a \textbf{mail server} delivers the messages.
  \item Using online banking: a \textbf{database server} stores and retrieves the account.
\end{itemize}

\field{What a server does:}{Stores files, hosts websites, manages databases, runs
applications, sends and receives email, shares printers and other resources, and manages
users and security.}

\subsection{Types of servers}

\begin{tabularx}{\textwidth}{@{}lX@{}}
\toprule
\rowcolor{tablehead} \textbf{Server type} & \textbf{Purpose} \\
\midrule
Web server & Hosts websites \\
File server & Stores and shares files \\
Database server & Stores databases \\
Mail server & Sends and receives emails \\
Application server & Runs business applications \\
Print server & Manages printers \\
DNS server & Converts domain names into IP addresses \\
\bottomrule
\end{tabularx}

\section{The Traditional Server Model}

\field{Short description:}{Before virtualization and cloud computing, organizations
typically deployed \textbf{one physical server for one application}. Each application had
its own dedicated hardware.}

\begin{tabularx}{\textwidth}{@{}lX@{}}
\toprule
\rowcolor{tablehead} \textbf{Application} & \textbf{Physical server} \\
\midrule
Website & Server 1 \\
Database & Server 2 \\
Email & Server 3 \\
ERP & Server 4 \\
\bottomrule
\end{tabularx}

\subsection{Advantages}

\begin{itemize}\itemsep3pt
  \item \textbf{High performance.} Applications have direct access to hardware resources.
  \item \textbf{Better security.} Each application runs on its own physical machine.
  \item \textbf{Reliability.} No interference from other applications sharing the hardware.
  \item \textbf{Simple architecture.} Easy to install, configure and troubleshoot.
  \item \textbf{Full hardware control.} Administrators control the server hardware completely.
\end{itemize}

\subsection{Disadvantages}

\begin{itemize}\itemsep3pt
  \item \textbf{Low resource utilization.} Most servers use only 10--20\% of CPU and memory,
        leaving the rest idle.
  \item \textbf{High cost.} Each application requires separate hardware.
  \item \textbf{High power consumption.} More servers draw more electricity and need more
        cooling.
  \item \textbf{More physical space.} Multiple servers occupy valuable rack space.
  \item \textbf{Difficult scalability.} A new application often means buying new hardware.
  \item \textbf{Higher maintenance.} Every server needs updates, monitoring, backups and
        repairs.
  \item \textbf{Single point of failure.} If a physical server fails, its application becomes
        unavailable.
\end{itemize}

\subsection*{Why a new technology was needed}

As organizations expanded they faced too many physical servers, rising hardware costs,
increased electricity and cooling expenses, underutilized resources, complex management, and
slow deployment of new services. These limitations led to \textbf{virtualization}, which
allows multiple virtual servers to run on a single physical machine, and virtualization in
turn became the foundation for \textbf{cloud computing}.

\newpage
\section{Introduction to Virtualization}

\field{Short description:}{Virtualization is a methodology for dividing computer resources
into more than one execution environment, applying concepts such as partitioning,
time-sharing, machine simulation and emulation.}

\field{Theory:}{Virtualization is a method in which multiple independent operating systems
run on one physical computer. It maximizes the usage of available physical resources, so one
can achieve high server usage: it increases total computing power while decreasing overhead.
Increasing the number of logical operating systems on a single host reduces hardware
acquisition and maintenance cost for an organization.}

\field{How it works:}{The architecture has three layers. Layer 1 is the physical
infrastructure of servers, networks and storage. Layer 2 is the virtual infrastructure,
consisting of hypervisors, virtual infrastructure services and automated solutions that
optimize the IT process. Layer 3 holds the virtual machines, where different operating
systems and applications are deployed. A single virtual infrastructure supports more than one
virtual machine, so the physical resources of the whole infrastructure are shared across many
virtual machines for maximum efficiency.}

\subsection{Architecture diagram}
\mmd[0.62]{01-virtualization-layers.png}

\subsection*{Notes}

\begin{itemize}\itemsep2pt
  \item Virtualization is a trend in IT that includes autonomic computing and utility
        computing: the environment manages itself, and every resource is seen as a utility
        that a client pays per use.
  \item It reduces the burden of workloads on users by centralizing administrative tasks and
        improving scalability.
\end{itemize}

\subsection{Reasons for using virtualization}

\begin{itemize}\itemsep3pt
  \item Virtual machines consolidate the workloads of under-utilized servers, saving on
        hardware, environmental costs and management.
  \item Legacy applications can keep running inside a VM.
  \item A VM provides a secure sandbox for running an untrusted application.
  \item A VM helps in building a secured computing platform.
  \item A VM provides an illusion of hardware.
  \item A VM simulates networks of independent computers.
  \item A VM supports running distinct operating systems with different versions.
\end{itemize}

\newpage
\section{Hypervisor}

\field{Short description:}{A hypervisor, or Virtual Machine Monitor (VMM), is software that
lets multiple operating systems run on a single physical machine.}

\field{Theory:}{It manages hardware resources such as CPU, memory and storage and allocates
them to virtual machines without interference. This improves hardware utilization, reduces
costs, and provides flexibility in cloud and server environments.}

\field{How it works:}{A hypervisor runs on hardware or on a host OS to create and manage
virtual machines, each with its own virtual CPU, memory, storage and network. It intercepts
guest OS requests and translates them to physical hardware, ensuring isolation, security and
stability.}

\subsection{Architecture diagram}
\mmd[0.34]{02-hypervisor-how-it-works.png}

\subsection{Type 1 hypervisor}

A Type 1 hypervisor runs directly on the host's hardware and does not rely on a host
operating system. This architecture offers better performance and security because there is
no intermediary OS. It is the standard for enterprise-level data centres and for cloud
providers such as Amazon Web Services and Microsoft Azure.

\field{Examples:}{VMware ESXi, Microsoft Hyper-V, KVM (Kernel-based Virtual Machine), Xen.}

\begin{itemize}\itemsep2pt
  \item \textbf{Pros:} high performance through direct hardware access; strong security with
        no intermediate OS layer; suitable for mission-critical workloads.
  \item \textbf{Cons:} requires dedicated hardware; setup and management are complex compared
        to Type 2.
\end{itemize}

\subsection{Type 2 hypervisor}

A Type 2 hypervisor runs on top of a conventional operating system such as Windows, macOS or
Linux. It is essentially an application within the host OS. This type is generally used for
desktop virtualization, development and testing, where a user needs to run multiple operating
systems on a personal computer. Performance is slightly lower than Type 1 because of the
overhead of the host OS.

\field{Examples:}{Oracle VM VirtualBox, VMware Workstation, Parallels Desktop.}

\begin{itemize}\itemsep2pt
  \item \textbf{Pros:} easy to install and use; useful for development, testing and malware
        analysis; good host--guest integration features.
  \item \textbf{Cons:} slower performance with no direct hardware access; security depends on
        the host OS, so a compromise of the host may affect the guests.
\end{itemize}

\subsection{Type 1 versus Type 2}
\mmd[0.44]{03-type1-vs-type2.png}

\newpage
\section{Full Virtualization}

\field{Short description:}{Full virtualization was introduced by IBM in 1966. It was the
first software solution for server virtualization and uses binary translation together with
direct execution.}

\field{How it works:}{The virtual machine completely isolates the guest OS from the
virtualization layer and the hardware. Kernel code is translated to replace non-virtualizable
instructions with new sequences of instructions that have the intended effect on the virtual
hardware, while user level code is executed directly on the processor for high performance.}

\field{Examples:}{Microsoft and Parallels systems.}

\subsection{Architecture diagram}
\mmd[0.30]{04-full-virtualization.png}

\section{Paravirtualization}

\field{Short description:}{Paravirtualization is the category of CPU virtualization that uses
hypercalls for operations handled at compile time.}

\field{How it works:}{The guest OS is not completely isolated, only partially isolated by the
virtual machine from the virtualization layer and hardware. The OS kernel is modified to
replace non-virtualizable instructions with hypercalls that communicate directly with the
virtualization layer, which improves performance and efficiency.}

\field{Examples:}{VMware and Xen.}

\subsection{Architecture diagram}
\mmd[0.30]{05-paravirtualization.png}

\newpage
\section{Hardware Based Virtualization}

\field{Short description:}{Hardware vendors added CPU features that simplify virtualization
directly in silicon.}

\field{How it works:}{First generation enhancements include Intel Virtualization Technology
(VT-x) and AMD's AMD-V. Both target privileged instructions with a new CPU execution mode that
allows the VMM to run in a root mode below ring 0. Privileged and sensitive calls are set to
trap automatically to the hypervisor, removing the need for either binary translation or
paravirtualization. The guest state is stored in Virtual Machine Control Structures (VT-x) or
Virtual Machine Control Blocks (AMD-V).}

\subsection{Architecture diagram}
\mmd[0.36]{06-hardware-assisted.png}

\newpage
\section{Types of Virtualization}

\subsection{Overview}
\mmd[0.24]{07-types-of-virtualization.png}

\subsection{Desktop virtualization}

Also called client virtualization. The virtualization happens on the user's site: desktops
are replaced with thin clients that use datacentre resources.

\subsection{Network virtualization}

Part of the virtualization infrastructure, used especially when virtualizing servers. It
allows multiple switches, VLANs and NAT-ing to be created in software.

\mmd[0.46]{10-network-virtualization.png}

\newpage
\subsection{Server and machine virtualization}

Virtualizing the server infrastructure, so no further physical servers are needed for
different purposes.

\subsection{Storage virtualization}

Widely used in datacentres with large storage. It helps create, delete and allocate storage
to different hardware, over a network connection. The leading technology here is the SAN.

\mmd[0.40]{11-storage-virtualization.png}

\subsection{Operating system level virtualization}

An abstraction layer between the traditional OS and user applications. OS-level virtualization
creates isolated \textbf{containers} on a single physical server, using one OS instance to
utilize the hardware and software in data centres. The containers behave like real servers. It
is commonly used to create virtual hosting environments that allocate hardware resources among
a large number of mutually distrusting users, and to a lesser extent to consolidate server
hardware by moving services on separate hosts into containers or VMs on one server.

\subsection{Application virtualization}

The technology isolates applications from the underlying operating system and from other
applications, to increase compatibility and manageability. Docker can be used for this
purpose.

\newpage
\section{Levels of Virtualization}

\field{Short description:}{A traditional computer runs a host operating system tailored to its
hardware architecture. After virtualization, different user applications managed by their own
guest operating systems can run on the same hardware, independent of the host OS. This is done
by adding a software layer called the virtualization layer, known as the hypervisor or virtual
machine monitor.}

\subsection{Before and after virtualization}
\mmd[0.42]{08-before-after-virtualization.png}

\subsection{Implementation levels}

Virtualization can be implemented at different levels within a computing system, offering
varying degrees of isolation and flexibility.

\mmd[0.20]{09-levels-of-virtualization.png}

\newpage
\subsection{Instruction set architecture level}

\begin{itemize}\itemsep2pt
  \item Virtualization is performed by emulating a given ISA with the ISA of the host machine.
  \item For example, MIPS binary code can run on an x86-based host machine through ISA
        emulation.
  \item It becomes possible to run a large amount of legacy binary code written for various
        processors on any given new hardware host.
\end{itemize}

\subsection{Hardware abstraction level}

\begin{itemize}\itemsep2pt
  \item Hardware-level virtualization is performed right on top of the bare hardware. It
        generates a virtual hardware environment for a VM, and the process manages the
        underlying hardware through virtualization.
  \item The idea is to virtualize a computer's resources, such as its processors, memory and
        I/O devices, to upgrade the hardware utilization rate by multiple concurrent users.
  \item It was implemented in the IBM VM/370 in the 1960s. More recently the Xen hypervisor
        has been applied to virtualize x86-based machines running Linux or other guest OS
        applications.
\end{itemize}

\subsection{Operating system level}

\begin{itemize}\itemsep2pt
  \item An abstraction layer between the traditional OS and user applications.
  \item It creates isolated containers on a single physical server, with the OS instances
        utilizing the hardware and software in data centres.
  \item The containers behave like real servers, and are commonly used for virtual hosting
        environments allocating hardware among mutually distrusting users.
\end{itemize}

\subsection{Library support level}

\begin{itemize}\itemsep2pt
  \item Most applications use APIs exported by user-level libraries rather than lengthy system
        calls, so a well-documented API becomes another candidate for virtualization.
  \item Virtualization with library interfaces controls the communication link between
        applications and the rest of the system through API hooks.
  \item WINE implements this approach to support Windows applications on top of UNIX hosts.
  \item vCUDA allows applications executing within VMs to leverage GPU hardware acceleration.
\end{itemize}

\subsection{User-application level}

\begin{itemize}\itemsep2pt
  \item Virtualization at the application level virtualizes an application as a VM. On a
        traditional OS an application runs as a process, so this is also known as
        process-level virtualization.
  \item The most popular approach is to deploy high level language (HLL) VMs. The
        virtualization layer sits as an application program on top of the operating system and
        exports the abstraction of a VM that runs programs compiled to a particular abstract
        machine definition.
  \item The Microsoft .NET CLR and the Java Virtual Machine are two good examples.
\end{itemize}

\newpage
\section{Containers}

\field{Short description:}{A container is an isolated, stand-alone unit that encapsulates an
application and all its dependencies. It runs consistently in any environment, independent of
the host system.}

\field{Theory:}{It is a light, executable software package that wraps everything an
application needs: libraries, configuration files and binaries. The application behaves the
same way on a developer's laptop, in a testing environment, or on a production server. Unlike
traditional virtual machines that carry a full OS, containers pack only what is required,
which makes them faster and more efficient.}

\field{How it works:}{Containerization is the process of packing an application together with
all its dependencies into a container so it runs consistently from one computing environment
to another. In simple terms it uses the host OS kernel to run many isolated instances of
applications on the same machine, which makes it very lightweight and efficient.}

\subsection{Containers versus virtual machines}

Both containers and VMs run applications in isolated environments, but they differ in
architecture, performance, resource usage and startup time.

\mmd[0.40]{12-containers-vs-vms.png}

\newpage
\subsection{Comparison}

\begin{tabularx}{\textwidth}{@{}lXX@{}}
\toprule
\rowcolor{tablehead} & \textbf{Containers} & \textbf{Virtual machines} \\
\midrule
Architecture & All containers share the host OS kernel; the running user spaces are isolated,
which makes them lightweight. & A hypervisor running on the host OS includes a full guest OS
with virtualized hardware. \\
\addlinespace
Boot time & Much less, typically seconds, as no full OS has to boot. & Longer, because the
full OS in the VM must be initialized. \\
\addlinespace
Isolation & At the process level, less strong than a VM, though for many use cases this does
not matter. & Very good, because each VM is a system on its own with its own OS. \\
\addlinespace
Resource usage & Fewer resources: only the necessary binaries and libraries, not an entire
OS. & Very high, as the full OS overhead is incurred for each instance. \\
\bottomrule
\end{tabularx}

\section{Container Orchestration}

\field{Short description:}{Container orchestration is the process of automating the
deployment, management, scaling and networking of containers throughout their lifecycle.}

\field{Theory:}{Orchestration platforms such as Kubernetes create a smooth, self-managing
operation and coordinate microservices held in separate containers. Kubernetes eliminates many
of the manual processes involved in deploying and scaling containerized applications, and Red
Hat creates essential tools for securing, simplifying and automatically updating container
infrastructure.}

\subsection{Why it is needed}

When many containers run across different machines, managing them manually becomes difficult.
An orchestration platform automates:

\begin{itemize}\itemsep3pt
  \item \textbf{Deployment:} starts and distributes containers across servers.
  \item \textbf{Scaling:} increases or decreases the number of container instances based on
        demand.
  \item \textbf{Load balancing:} distributes incoming traffic among healthy containers.
  \item \textbf{Self-healing:} restarts failed containers or moves them to healthy nodes.
  \item \textbf{Service discovery:} allows containers to find and communicate with each other.
  \item \textbf{Rolling updates and rollbacks:} updates applications with minimal downtime and
        reverts changes if needed.
  \item \textbf{Resource management:} efficiently allocates CPU, memory and storage.
\end{itemize}

\subsection*{Benefits}

Orchestration platforms automate the entire lifecycle of containers, transforming what would
be a non-stop maintenance operation for IT teams into a smooth, self-managing one, and bring
advantages in development methods, costs and security.

\end{document}
